\section*{الفصل \ref{ch:b2:batteries-electrolysis} --- البطاريات وخلايا الوقود والتحليل الكهربائي}

\begin{solution}{exo:b2:batteries-electrolysis:1}
الزنك وحده: نحو \qty{-0.82}{V} وتيار قدره 0.15 (بوحدات النموذج). والزنك الملامس
للنحاس: نحو \qty{-0.74}{V} و2.4. فالتيار، ومن ثم سرعة الذوبان،
أكبر بنحو ست عشرة مرة عند التماس مع النحاس.
\end{solution}

\begin{solution}{exo:b2:batteries-electrolysis:2}
$m = 500 \times 3600 \times 65.38/(2 \times 96\,485) = \qty{610}{g}$.
\end{solution}

\begin{solution}{exo:b2:batteries-electrolysis:3}
$E^\circ = \qty{2.04}{V}$ (كما في الفصل)؛ وست خلايا على التوالي تعطي
\qty{12}{V}.
\end{solution}

\begin{solution}{exo:b2:batteries-electrolysis:4}
$1/6.94 = \qty{0.144}{mol}$؛ $0.144 \times 96\,485 = \qty{1.39e4}{C}$، أي
\qty{3.86}{A.h}.
\end{solution}

\begin{solution}{exo:b2:batteries-electrolysis:5}
بمردود 100~\%: $200 \times 86\,400 \times 63.546/(2 \times 96\,485) = \qty{5.69}{kg}$؛
$\eta_F = 5.50/5.69 = 0.967$.
\end{solution}

\begin{solution}{exo:b2:batteries-electrolysis:6}
$w = 2 \times 96\,485 \times 3.3/(0.90 \times 0.06538) = \qty{1.08e7}{J/kg}$، أي
\qty{3.0}{kW.h/kg}.
\end{solution}

\begin{solution}{exo:b2:batteries-electrolysis:7}
$U = E - rI$: $r = (1.55 - 1.40)/(0.60 - 0.10) = \qty{0.30}{\ohm}$؛ $E = 1.55 + 0.30
\times 0.10 = \qty{1.58}{V}$.
\end{solution}

\begin{solution}{exo:b2:batteries-electrolysis:8}
المصعد \ce{2Cl- -> Cl2 + 2e-}؛ والمهبط \ce{2H2O + 2e- -> H2 + 2OH-}. $\Delta_r G^\circ =
2(-157.24) - 2(-131.23) - 2(-237.13) = \qty{422.2}{kJ/mol}$، $U_{\min} =
\qty{2.19}{V}$. فالكلور كان سيتفاعل مع الهيدروكسيد (\ce{Cl2 + 2OH- -> Cl- +
ClO- + H2O})، فيُفسد الناتجين كليهما، وخليط الكلور والهيدروجين
قد ينفجر.
\end{solution}

\begin{solution}{exo:b2:batteries-electrolysis:9}
\ce{NaCl -> Na + 1/2Cl2}، $n = 1$: $U_{\min} = 310\,674/96\,485 = \qty{3.22}{V}$.
ولكل كيلوغرام من الصوديوم: $96\,485 \times 3.22/0.02299 = \qty{1.35e7}{J}$، أي
\qty{3.75}{kW.h}.
\end{solution}

\begin{solution}{exo:b2:batteries-electrolysis:10}
على مهبط من الزنك يكون اختزال \ce{H+} بطيئًا ويقع جداره تحت
موجة الزنك: فكلما خُفض الجهد ترسب الزنك أولًا، مع قليل من
الهيدروجين. أما على البلاتين فاختزال \ce{H+} سريع ويبدأ قرب
\qty{0}{V}: فلا يتصاعد في البداية إلا الهيدروجين. وبمجرد أن يُغطى البلاتين
بالزنك، يسلك المهبط سلوك مهبط من الزنك فيترسب الزنك.
\end{solution}

\begin{solution}{exo:b2:batteries-electrolysis:11}
عند \qty{1250}{K}: $\Delta_f G^\circ(\ce{Al2O3}) = -1328.3 + 0.75 \times 66.0 =
\qty{-1278.8}{kJ/mol}$، $\Delta_f G^\circ(\ce{CO2}) = \qty{-396.1}{kJ/mol}$. $\Delta_r G^\circ =
3(-396.1) - 2(-1278.8) = \qty{1369}{kJ}$، $n = 12$: $U_{\min} = \qty{1.18}{V}$. ومع مصعد خامل:
$\Delta_r G^\circ = \qty{2557.5}{kJ}$، $U_{\min} = \qty{2.21}{V}$؛ فالكربون المحترق
يوفر نحو \qty{1.03}{V} من الشغل.
\end{solution}

\begin{solution}{exo:b2:batteries-electrolysis:12}
لكل مول من الماء تكون الشحنة $2F$ نفسها في الاتجاهين: المردود $0.70/1.80 =
0.39$. والباقي حرارة: \omterm{def:b2:current-potential-curves:overpotential}{فرط جهد} قطبي الأكسجين كليهما (البطيئين في
الاتجاهين)، \omterm{def:b2:current-potential-curves:overpotential}{وفرط جهد} قطبي الهيدروجين، والخسائر الأومية.
\end{solution}

\begin{solution}{pb:b2:batteries-electrolysis:1}
\textbf{1.} الألومنيوم $+3 \to 0$؛ ويبقى الأكسجين $-2$؛ والكربون $0 \to +4$.
\textbf{2.} المهبط \ce{Al^3+ + 3e- -> Al}؛ والمصعد \ce{C + 2O^2- -> CO2 + 4e-}.
\textbf{3.} \ce{2Al2O3 + 3C -> 4Al + 3CO2}، $n = 12$.
\textbf{4.} \ce{Al2O3}: $-1328.3 + 0.75 \times (1328.3 - 1262.3) = \qty{-1278.8}{kJ/mol}$؛
\ce{CO2}: \qty{-396.1}{kJ/mol}.
\textbf{5.} $3(-396.1) - 2(-1278.8) = \qty{+1369}{kJ}$.
\textbf{6.} $1\,369\,000/(12 \times 96\,485) = \qty{1.18}{V}$.
\textbf{7.} $2 \times 1278.8/(12 \times 96.485) = \qty{2.21}{V}$: فأكسدة الكربون،
الطاردة للطاقة بذاتها، تدفع نحو نصف شغل تفكيك الأكسيد.
\textbf{8.} $10^6/26.982 = \qty{3.706e4}{mol}$.
\textbf{9.} $3 \times 96\,485 \times 3.706 \times 10^4 = \qty{1.073e10}{C}$.
\textbf{10.} $1.073 \times 10^{10}/0.94 = \qty{1.141e10}{C}$.
\textbf{11.} $1.141 \times 10^{10}/3.0 \times 10^5 = \qty{3.80e4}{s}$، أي \qty{10.6}{h}.
\textbf{12.} $24/10.6 = 2.27$: أي \qty{2.27}{t} في اليوم.
\textbf{13.} جزء من الألومنيوم المذاب في الحوض يبلغ غاز المصعد
ويعيد ثنائي أكسيد الكربون أكسدته، فتُنفق شحنته مرتين؛ ويتسرب أيضًا بعض
التيار عبر الحوض دون أن يحلله كهربائيًا.
\textbf{14.} $1.141 \times 10^{10} \times 4.1 = \qty{4.68e10}{J}$، أي \qty{13.0}{MW.h}.
\textbf{15.} \qty{13.0}{kW.h/kg}.
\textbf{16.} $3 \times 96\,485 \times 1.18/0.026982 = \qty{1.27e7}{J/kg}$، أي \qty{3.5}{kW.h/kg}.
\textbf{17.} إلى حرارة: \omterm{def:b2:current-potential-curves:overpotential}{فرط جهد} القطبين، وأكثرها الهبوط
الأومي في الحوض. وهذه الحرارة تُبقي الحوض منصهرًا عند نحو
\qty{1250}{K}، وهو ما كان سيحتاج وقودًا لولاها.
\textbf{18.} $4.1 \times 3.0 \times 10^5 = \qty{1.23}{MW}$.
\textbf{19.} $28\,690/0.026982 = \qty{1.06e6}{J/kg}$، أي \qty{0.30}{kW.h/kg}: نحو 2~\% من
طاقة \omterm{def:b2:batteries-electrolysis:electrolysis}{التحليل الكهربائي}.
\textbf{20.} $1.18/4.1 = 0.29$.
\textbf{21.} $3.706 \times 10^4 \times 3/4 \times 12.011 = \qty{334}{kg}$.
\textbf{22.} $3.706 \times 10^4 \times 3/4 \times 44.009 = \qty{1.22e3}{kg}$.
\textbf{23.} تحترق الأطراف الساخنة للمصاعد في الهواء؛ ويغادر جزء من الكربون
على هيئة أحادي أكسيد الكربون (\ce{C + CO2 -> 2CO} عند تلك الدرجة)، الذي يأخذ ضعف
كمية الكربون لكل ذرة أكسجين.
\textbf{24.} \qty{1.22}{kg} من \ce{CO2} لكل كيلوغرام من الألومنيوم.
\textbf{25.} كان الجهد الأدنى سيرتفع بنحو \qty{1.03}{V}، لكن المصعد
كان سيطلق الأكسجين بدل ثنائي أكسيد الكربون ولا يُستهلك.
\textbf{26.} تستعمل الخلية $\boldsymbol{\approx \qty{13}{kW.h}}$ من الطاقة
الكهربائية لكل كيلوغرام من الألومنيوم.
\end{solution}
